\documentclass[12pt]{article}
\usepackage[T1]{fontenc}
\usepackage[utf8]{inputenc}
\usepackage{lmodern}
\usepackage{textcomp}
\usepackage{enumitem}
\usepackage{geometry}
\geometry{margin=1in}
\begin{document}
\begin{center}
Answer Key - Chemistry - Undergraduate Level Assessment 16 \
\small (Answers are ordered exactly as in the question paper.)
\end{center}

\section*{Multiple Choice}
\begin{enumerate}[label=\arabic*.]
\item \textbf{Answer:} C
\\ \textit{Options:} A) 54.91 MHz; B) 239.2 MHz; C) 345.0 MHz; D) 2167 MHz
\\ \textit{Note:} The NMR frequency of 31 P in a 20.0 T magnetic field is calculated using the formula ν = γB 0/2π, where γ (the gyromagnetic ratio for 31 P) is approximately 40.5 MHz/T, resulting in a frequency of about 345.0 MHz when substituting the values.
\item \textbf{Answer:} D
\\ \textit{Options:} A) The most common oxidation state for the lanthanide elements is +3.; B) Lanthanide complexes often have high coordination numbers (\textgreater{} 6).; C) All of the lanthanide elements react with aqueous acid to liberate hydrogen.; D) The atomic radii of the lanthanide elements increase across the period from La to Lu.
\\ \textit{Note:} The atomic radii of the lanthanide elements actually decrease across the period from La to Lu due to the increasing effective nuclear charge which pulls the electrons closer to the nucleus, making the statement that they increase incorrect.
\item \textbf{Answer:} D
\\ \textit{Options:} A) Both 35 Cl and 37 Cl have I = 0.; B) The hydrogen atom undergoes rapid intermolecular exchange.; C) The molecule is not rigid.; D) Both 35 Cl and 37 Cl have electric quadrupole moments.
\\ \textit{Note:} The 1 H NMR spectrum of 12 CHCl 3 appears as a singlet because the presence of 35 Cl and 37 Cl isotopes, which have electric quadrupole moments, leads to averaging of the magnetic interactions, resulting in no splitting of the hydrogen signal.
\item \textbf{Answer:} D
\\ \textit{Options:} A) The process is exothermic.; B) The process does not involve any work.; C) The entropy of the system increases.; D) The total entropy of the system plus surroundings increases.
\\ \textit{Note:} A spontaneous process is characterized by an increase in the total entropy of the system and its surroundings, in accordance with the second law of thermodynamics, indicating that natural processes tend to move towards a state of greater disorder or randomness.
\item \textbf{Answer:} D
\\ \textit{Options:} A) 10.8 x 10^-5; B) 4.11 x 10^-5; C) 3.43 x 10^-5; D) 1.71 x 10^-5
\\ \textit{Note:} The correct answer of 1.71 x 10^-5 is derived from the Boltzmann distribution, which calculates the equilibrium polarization of nuclear spins in a magnetic field by considering the ratio of population differences between energy states at a given temperature, specifically for the spin-1/2 13 C nuclei in a 20.0 T magnetic field.
\end{enumerate}

\section*{True / False}
\begin{enumerate}[label=\arabic*.]
\setcounter{enumi}{5}
\item \textbf{Answer:} False
\\ \textit{Options:} A) True; B) False
\\ \textit{Note:} The thermal stability of group-14 hydrides actually decreases down the group, with CH 4 being the most stable and PbH 4 the least stable, contrary to the stated order.
\item \textbf{Answer:} False
\\ \textit{Options:} A) True; B) False
\\ \textit{Note:} The orbital angular momentum quantum number, l, for the electron most easily removed from ground-state aluminum is 1 (corresponding to the 3 p orbital), not 2, which is associated with the 3 d orbital.
\item \textbf{Answer:} True
\\ \textit{Options:} A) True; B) False
\\ \textit{Note:} The +1 oxidation state of thallium is more stable than the +3 oxidation state due to the presence of a filled d-subshell, which provides greater stability and reduces the effective nuclear charge experienced by the outer electrons.
\item \textbf{Answer:} True
\\ \textit{Options:} A) True; B) False
\\ \textit{Note:} The magnetogyric ratio of 23 Na is indeed 7.081 x $10^7$ T^-1 s^-1, as this value is a well-established constant in nuclear magnetic resonance (NMR) spectroscopy specific to the sodium-23 isotope.
\item \textbf{Answer:} True
\\ \textit{Options:} A) True; B) False
\\ \textit{Note:} The difference in infrared absorption frequency between DCl and HCl is primarily due to their reduced mass because the reduced mass influences the vibrational frequency of the bond, with DCl having a greater reduced mass than HCl, resulting in lower vibrational frequency and thus a shift in absorption.
\end{enumerate}

\section*{Long Form Response}
\begin{enumerate}[label=\arabic*.]
\setcounter{enumi}{10}
\item \textbf{Answer:} Metal ions play a significant role as paramagnetic quenchers in chemical reactions by interacting with excited states of molecules, leading to non-radiative decay of energy. Paramagnetic ions, such as $Mn^2$⁺ and $Fe^3$⁺, can effectively reduce fluorescence and phosphorescence by facilitating spin-forbidden transitions. However, $Cr^3$⁺ is not suitable for this purpose because it has a low magnetic moment and is typically found in a diamagnetic state in its most common electronic configuration. This lack of unpaired electrons limits its ability to act as an effective quencher, as it cannot significantly interact with the excited state of the molecules involved. Consequently, while other metal ions can enhance quenching through their paramagnetic properties, $Cr^3$⁺ fails to provide the necessary magnetic characteristics for such applications.
\\ \textit{Note:} The answer correctly identifies that metal ions like Mn²⁺ and Fe³⁺ are effective paramagnetic quenchers due to their ability to promote non-radiative energy decay through spin-forbidden transitions, while Cr³⁺ is unsuitable because it possesses a low magnetic moment and often exists in a diamagnetic state, limiting its quenching capability.
\item \textbf{Answer:} The principal allotropes of elemental boron include amorphous boron, α-rhombohedral boron, and β-rhombohedral boron, each displaying unique solid-state structures. Central to the stability and formation of these allotropes are the B$_{12}$ icosahedra, which consist of twelve boron atoms arranged in a highly symmetrical, closed structure. These icosahedra serve as fundamental building blocks, enabling the formation of a three-dimensional network that contributes to the overall stability of the allotropes. In α-rhombohedral boron, for instance, the arrangement of B$_{12}$ icosahedra leads to a stable crystal lattice, while in β-rhombohedral boron, the interconnection of these icosahedra results in a denser packing and enhanced stability. Thus, the presence and arrangement of B$_{12}$ icosahedra are crucial in determining the structural properties and stability of boron's various allotropes.
\\ \textit{Note:} The answer correctly identifies the principal allotropes of elemental boron and emphasizes the significance of B$_{12}$ icosahedra as essential structural units that provide stability and enable the formation of diverse solid-state configurations.
\end{enumerate}

\vfill
\noindent\textit{End of Answer Key}
\end{document}